% E.30 onward reserved for the five-point census.
\subsection{Five marked points and cote 148}
\label{sec:E2}

We distinguish three classification problems.  The first concerns
a complex sphere with a five-point set $K$.  The second equips
that configuration with an antiholomorphic involution.  The third
concerns finite subgroups of the extended mapping class group of
the five-punctured sphere.  These are different problems.
An antiholomorphic map of finite order can have order greater than two.
We write $D_n$ for the
dihedral group of order $2n$ and $\mu_n$ for the complex $n$th roots
of unity.  The manuscript passages here are cote 148, archival
pp.~30--35, scans 31--36; its separate classification occupies
handwritten pp.~1--3, scans 34--36.  We record below which conclusions
come from independent mathematics and which diagram entries require
interpretation of the manuscript's additional decorations.  We inspected
the original Montpellier PDF directly \cite{DE_Grothendieck148}.  Its printed page numbers are
one less than the one-based PDF scan numbers; the supplied first-pass
transcription shifts the archival page numbers down by one further unit.

\statement{E.30}{Proved}{The holomorphic classification of five-point configurations}
The possible full holomorphic automorphism groups are
\[
 1,\quad C_2,\quad C_4,\quad D_3,\quad D_5.
\]
If the maximum order $\nu$ of a nonidentity holomorphic automorphism
is $5$, $4$, or $3$, there is exactly one isomorphism class, respectively
\[
 \mu_5,\qquad\{\infty\}\cup\mu_4,
 \qquad\{0,\infty\}\cup\mu_3.
\]
Their full groups are $D_5,C_4,D_3$.  The labeled structures are,
respectively, an unoriented pentagon, a distinguished point with an
oriented square on the other four points, and a three-set together
with its two-element orientation torsor.  These assertions verify
the holomorphic classification on cote 148, scan 34, handwritten p.~1.
\begin{proof}
The holomorphic group acts faithfully on $K$, since a M\"obius
transformation fixing three points is the identity.  It is finite.
A nonidentity finite-order M\"obius map is conjugate to
$z\mapsto\zeta_\nu z$.  Its only fixed points are $0,\infty$, and
all other orbits have length $\nu$.  Hence
$5=f+m\nu$, with $0\leq f\leq2$, $m\geq1$, giving
$\nu\in\{2,3,4,5\}$ and the three displayed normal forms for
$\nu\geq3$.

For $K=\mu_5$, every automorphism preserves its unique generalized
circle and hence its cyclic order up to reversal.  It acts on the
five points through $D_5$.  The maps $z\mapsto\zeta_5z$ and
$z\mapsto1/z$ realize all of $D_5$.
For $K=\{\infty\}\cup\mu_4$, the unit circle is the unique
generalized circle containing four marked points: any other
four-set contains $\infty$, and its other three points are not
collinear.  Thus $\infty$ is invariant.  An automorphism is affine;
preservation of the centroid of $\mu_4$ forces its translation
part to vanish, leaving exactly multiplication by $\mu_4$.
For $K=\{0,\infty\}\cup\mu_3$, deleting either pole gives an
equianharmonic four-set, with $j=0$, whereas deleting a root gives
$j=2048/3$.  This follows by substitution in
\[
 j(\lambda)=256\frac{(1-\lambda+\lambda^2)^3}
 {\lambda^2(1-\lambda)^2}.
\]
Therefore the pole pair is intrinsic.  An automorphism has form
$az$ or $a/z$, and preservation of $\mu_3$ forces $a\in\mu_3$.
The resulting group is $D_3$.

If every nonidentity element has order two, the group is elementary
abelian.  After one involution is written as $z\mapsto-z$, every
commuting involution is either that map or $z\mapsto a/z$.
There are at most four elements.  A Klein four action has no common
fixed point and all its point orbits have even cardinality, so cannot
preserve a five-set.  Thus the remaining nontrivial full group is
$C_2$.  The stated labeled structures have stabilizers respectively
$D_5,C_4,S_3$ in $S_5$, with the same actions as the normal forms;
this also proves the categorical identifications used in scan 34.
\end{proof}

\statement{E.31}{Proved}{The one-parameter involution locus and the square and tetrahedron cases}
Every five-point configuration with a nontrivial holomorphic
involution can be written
\[
 K_t=\{\infty,1,-1,t,-t\},\qquad t\ne0,\ \ t^2\ne1.
\]
For a configuration with the involution specified, its isomorphism
class is determined by
\[
 \lambda=t^2,\qquad\beta=\lambda+\lambda^{-1}\in\mathbb C\setminus\{2\}.
\]
The full group is $C_2$ except at
\[
 \begin{array}{c|ccc}
 \beta&-2&-10/3&18\\\hline
 \operatorname{Aut}^+(K_t)&C_4&D_3&D_5.
 \end{array}
\]
A harmonic four-set together with a fixed point of one of its
fixed-point-free involutions gives $\beta=-2$ or $34$.
The second case has full group $C_2$; it is neither the pentagonal
nor the square-pyramidal case.  An equianharmonic four-set with
such a fixed point gives $\beta=-14$ and full group $C_2$.
These formulas complete the calculation begun and then crossed out
on cote 148, scan 35, handwritten p.~2, and answer its question
about adding an edge midpoint to a spherical square.
\begin{proof}
An involution preserving five marks has exactly one marked fixed
point.  Send it to $\infty$ and the other fixed point to $0$,
write the involution as $z\mapsto-z$, and scale one two-orbit to
$\{1,-1\}$.  Changing the chosen two-orbit replaces $t$ by
$\pm t^{-1}$; changing its ordering replaces $t$ by $-t$.
Conversely any isomorphism preserving the specified involution
must preserve its marked and unmarked fixed points individually,
so is a scaling.  These are exactly the identifications.  Hence
$\beta$ is complete and $\beta=2$ is exactly the forbidden collision.
The normal forms of E.30 give $t=i$, $i/\sqrt3$, and $\sqrt5-2$,
respectively, yielding the three exceptional values.  E.30 excludes
any other enhancement.  All involutions in each of $C_4,D_3,D_5$
are conjugate within that group, so forgetting the chosen involution
does not introduce another value of $\beta$.

For the four-set $\{1,-1,t,-t\}$ a cross-ratio is
$((1-t)/(1+t))^2$, and direct substitution gives
\[
 j=16\frac{(\lambda^2+14\lambda+1)^3}
 {\lambda(\lambda-1)^4}.
\]
The equianharmonic condition is $\lambda^2+14\lambda+1=0$,
namely $\beta=-14$.  The harmonic condition factors exactly as
\begin{align*}
 &16(\lambda^2+14\lambda+1)^3
 -1728\lambda(\lambda-1)^4\\
 &\hspace{15mm}=16(\lambda+1)^2(\lambda^2-34\lambda+1)^2,
\end{align*}
giving $\beta=-2$ or $34$.  In the square model $\mu_4$, the
involution $-z$ gives the pole and $\beta=-2$; an involution such
as $i/z$ gives a point halfway between consecutive vertices and
$\beta=34$.  For the tetrahedral four-set the three involutions
have six fixed points, the vertices of the associated octahedron;
adding any one gives $\beta=-14$.
The exact symbolic program \texttt{check\_special\_configurations.py}
enumerates every M\"obius candidate from its images on three marks
and verifies all six listed normal forms, including the generic
control $t=2$, with no numerical tolerance.
\end{proof}

\statement{E.32}{Refuted}{Every finite antiholomorphic automorphism preserving five marks has order two}
The inference in cote 148, scan 35, handwritten p.~2, that an
even-order antiholomorphic automorphism preserving five points
must fix a marked point and therefore have order two is false.
The exact possible orders of antiholomorphic automorphisms of
five-point configurations are $2,6,10$.
\begin{proof}
Let $\omega=e^{2\pi i/3}$ and
$K=\{0,\infty,1,\omega,\omega^2\}$.  The map
\[
 \sigma(z)=\frac{\omega}{\overline z}
\]
interchanges $0$ and $\infty$ and cyclically permutes the three
roots.  Its marked permutation has cycle type $(2)(3)$;
$\sigma^2(z)=\omega^2z$, so its exact order is six.
It has no fixed point: a finite nonzero fixed point would satisfy
$|z|^2=\omega$, which is impossible, and the two poles are exchanged.
For $K=\mu_5$ the map
\[
 \tau(z)=\frac{\zeta_5}{\overline z}
\]
has order ten.  Its restriction to $K$ is a five-cycle, while
$\tau^5(z)=1/\overline z$ fixes all five marks.  Thus restriction
to the marks alone is not a faithful representation of the extended
automorphism group.  The pair consisting of that permutation and
the orientation sign is faithful: an element in its kernel is a
holomorphic map fixing five points.

For completeness, any antiholomorphic automorphism has even order,
and its square is in one of the positive groups in E.30.  Orders
four or eight would force a marked fixed point, since a permutation
of five letters whose order divides four and which has no fixed
point cannot occur.  A finite-order anti-M\"obius map with a fixed
point is an involution: send the point to $\infty$, write the map
as $a\overline z+b$, and observe that its square has positive real
linear coefficient $|a|^2$, so finite order forces the square to
be the identity.  The remaining possibilities are $2,6,10$, all
realized above or by conjugation.  The finite semilinear and
permutation computations are checked in \texttt{check\_census.py}.
\end{proof}

\statement{E.33}{Proved}{Real involutions, dimensions, and their exact relation to the symmetric square}
An antiholomorphic \emph{involution} preserving five marked points
has a fixed circle and fixes exactly $r=1,3,$ or $5$ of the marks.
Each of these three spaces of real pairs $(\mathbb P^1,K,\sigma)$
has real dimension two.  The examples in cote 148, scan 36,
handwritten p.~3, are valid: the pentagonal configuration has an
$r=5$ real structure, the triangular bipyramid has an $r=3$
structure, and the square pyramid has an $r=1$ structure.
These examples are not exclusive to the indicated type.
\begin{proof}
An involution of an odd set has a fixed point.  Hence $\sigma$ has
a fixed point on the sphere and is M\"obius-conjugate to complex
conjugation.  The marks off its fixed circle occur in conjugate
pairs, giving $r=1,3,5$.  For $r=5$, five distinct real points
modulo $\operatorname{PGL}_2(\mathbb R)$ have dimension $5-3=2$.
For $r=3$, choose and order the real marks as $0,1,\infty$; then
\[
 K=\{0,1,\infty,\lambda,\overline\lambda\},
 \qquad\operatorname{Im}\lambda>0.
\]
The remaining finite changes of normalization do not alter dimension.
For $r=1$, place the real mark at $\infty$; two distinct points
in the upper half-plane and their conjugates give four real
parameters, modulo the two-dimensional real affine group.
Again the finite exchanges do not alter dimension.

The complete real-structure counts on the three enhanced
configurations are
\[
\begin{array}{c|c|c}
 K&\operatorname{Aut}^{\pm}(K)&
 \text{fixed-mark counts of its negative involutions}\\\hline
 \mu_5&D_5\times C_2&5,1,1,1,1,1\\
 \{0,\infty\}\cup\mu_3&D_3\times C_2&3,3,3,3\\
 \{\infty\}\cup\mu_4&D_4&3,3,1,1
\end{array}
\]
For example $z\mapsto1/\overline z$ fixes $\mu_5$, conjugation
fixes the two poles and $1$ in the bipyramid, and
$z\mapsto i\overline z$ fixes only $\infty$ in the square pyramid.
The other entries follow by multiplying these maps by the positive
groups of E.30 and counting their fixed marks.

One must also qualify the real-locus shorthand on scan 36(a).
After three ordered real marks are normalized, the model
$\operatorname{Sym}^2_{\mathrm{dist}}(
\mathbb P^1\setminus\{0,1,\infty\})$ parametrizes an unordered
remaining pair.  Its real points include both two real points and
a nonreal conjugate pair.  Thus its entire real locus is not the
$r=5$ locus: the pair $\{i,-i\}$ is an explicit $r=3$ counterexample
to that unqualified identification.  The $r=5$ locus is its
split-real part.  Fully unordered five-point moduli additionally
identifies different choices of the normalized triple.
\end{proof}

\statement{E.34}{Proved}{The generic real involution locus}
For $K_t$ of E.31 with full positive group $C_2$, a real structure
exists exactly when $\beta$ is real.  The full extended group is
then $C_2\times C_2$, and the fixed-mark counts of its two negative
involutions are
\[
\begin{array}{c|c}
 \text{range of }\beta&\text{counts}\\\hline
 \beta>2&\{5,1\}\\
 \beta<-2&\{3,3\}\\
 -2<\beta<2&\{1,1\}.
\end{array}
\]
The enhanced values in E.31 are omitted from the indicated ranges.
\begin{proof}
The invariant of the conjugate configuration is $\overline\beta$,
and E.31 proves completeness of $\beta$.  Thus a negative
isomorphism implies $\beta=\overline\beta$.  Conversely if
$\beta>2$, take $t$ real and use $\overline z$ and
$-\overline z$.  If $\beta<-2$, take $t$ purely imaginary and
use the same two maps.  If $-2<\beta<2$, choose $|t|=1$ and use
$t\overline z$ and $-t\overline z$.  Their respective fixed-mark
counts are exactly the displayed values.  An anti-isomorphism
normalizes the unique positive involution; it fixes its unique
marked pole and therefore has the form $a\overline z$ in this
normalization.  It is an involution.  The negative coset has two
elements, proving both exhaustiveness and the group assertion.
\end{proof}

\statement{E.35}{Proved}{Reflection section classes for five punctures}
Let $U=\mathbb P^1\setminus K$ and let $\sigma$ be a real involution
fixing $r\in\{1,3,5\}$ marks.  Put $m=(5-r)/2$.  Its discrete
orbifold extension is
\[
 1\longrightarrow\pi_1(U)\longrightarrow
 E_\sigma\longrightarrow C_2\longrightarrow1,
 \qquad
 E_\sigma\cong\Bigl(\mathop{*}_{j=1}^{r}C_2\Bigr)*F_m.
\]
The sign map is nontrivial on each $C_2$ factor and trivial on
$F_m$.  It has exactly $r$ conjugacy classes of sections under
$\pi_1(U)$, one for each component of $U^\sigma$; every section
has trivial fixed subgroup in $\pi_1(U)$.
This verifies the five-point instances of LM4 B.4 and B.5.
\begin{proof}
The quotient of the sphere by the reflection is a disk.  The $r$
real punctures split its mirror boundary into $r$ open intervals;
the $m$ conjugate pairs give $m$ interior punctures.  Cut from each
interior puncture to the boundary, and use paths from one interior
basepoint to the $r$ mirror intervals.  Orbifold van Kampen gives
one reflection generator of order two per interval and one free
orientation-preserving generator per interior puncture, with no
other relation.  The index-two kernel is torsion-free, and its rank
is $(r-1)+2m=4$, as required for a five-punctured sphere.

Reduced-word normal form in a free product shows that every
finite-order element is conjugate into one finite factor.  Distinct
factors have nonconjugate involutions, and the centralizer of a
factor involution is that factor itself.  Each $E_\sigma$-conjugacy
class of such an involution is already one $\pi_1(U)$-conjugacy
class: for a sign-negative conjugator $g$ of $t_j$, write
$g=(gt_j)t_j$ with $gt_j\in\pi_1(U)$.  Therefore there are exactly
$r$ section classes, and intersecting its centralizer with the
kernel gives $1$.  The quotient construction identifies the
classes with the $r$ open fixed intervals, giving the claimed
geometric correspondence.
\end{proof}

\statement{E.36}{Proved}{The twenty finite-subgroup classes in type $(0,5)$}
The discrete extended mapping class group of a sphere with five
unordered punctures has exactly twenty conjugacy classes of finite
subgroups, including the trivial group.  Seven preserve orientation
and thirteen have nontrivial sign.  In the following table, the
last column lists the fixed-mark counts of all negative involutions,
with multiplicity; an empty entry means there are none.
\begin{center}\small
\begin{tabular}{c|r|c|l}
 $G$&$|G|$&$G^+$&\text{negative involutions}\\\hline
 $1$&1&$1$&\\
 $C_2$&2&$C_2$&\\
 $C_3$&3&$C_3$&\\
 $C_4$&4&$C_4$&\\
 $C_5$&5&$C_5$&\\
 $D_3$&6&$D_3$&\\
 $D_5$&10&$D_5$&\\\hline
 $C_2$&2&$1$&$5$\\
 $C_2$&2&$1$&$3$\\
 $C_2$&2&$1$&$1$\\
 $C_2\times C_2$&4&$C_2$&$5,1$\\
 $C_2\times C_2$&4&$C_2$&$3,3$\\
 $C_2\times C_2$&4&$C_2$&$1,1$\\
 $C_6$&6&$C_3$&$3$\\
 $D_3$&6&$C_3$&$3,3,3$\\
 $D_4$&8&$C_4$&$1,1,3,3$\\
 $C_{10}$&10&$C_5$&$5$\\
 $D_5$&10&$C_5$&$1,1,1,1,1$\\
 $D_3\times C_2$&12&$D_3$&$3,3,3,3$\\
 $D_5\times C_2$&20&$D_5$&$1,1,1,1,1,5$
\end{tabular}
\end{center}
Thus the type $(0,5)$ total in LM4 B.26 is correct, although it
must not be identified with the number of full automorphism groups
of individual complex configurations.
\begin{proof}
Kerckhoff's realization theorem allows orientation-reversing mapping
classes, and its extension to punctures is explicit in
\cite[Section IV, pp.~255--256; Theorem 5]{DE_Kerckhoff1983}.
Hence a finite subgroup is realized by holomorphic and antiholomorphic
automorphisms of a five-punctured sphere, and extends over its
punctures.  E.30 shows that its positive subgroup is one of
$1,C_2,C_3,C_4,C_5,D_3,D_5$.
This agrees with the positive conjugacy classification in
\cite[Theorems 3--4]{DE_Stukow2004}.  We give the remaining
classification directly.

If $G^+=1$, a negative group is generated by a reflection, with
$r=1,3,5$, giving three classes.  If $G^+=C_2$, its unique marked
fixed pole is preserved by all of $G$.  In coordinates in which
$G^+=\langle-z\rangle$, any negative normalizer preserving that
pole has form $a\overline z$ and has order two.  The full group
is a Klein four group with two perpendicular reflection circles.
The other four marks are either one free four-orbit, two two-orbits
on one mirror, or one two-orbit on each mirror.  These give exactly
$\{1,1\},\{5,1\},\{3,3\}$.

For $G^+=C_3$, the marks are the two poles and one three-orbit.
A negative normalizer either fixes the poles separately, giving
the dihedral group generated by $\zeta_3z$ and $\overline z$,
or exchanges them, giving the cyclic group generated by
$\zeta_3/\overline z$.  For $G^+=C_4$, the single marked pole
must be fixed, and the only extension is the dihedral group
$D_4$.  For $G^+=C_5$, the same normalizer calculation gives
$D_5$ or $C_{10}$, according as the two poles are fixed separately
or exchanged.  Finally the $D_3$ and $D_5$ positive actions have
the rigid configurations of E.30; their full extended groups are
$D_3\times C_2$ and $D_5\times C_2$.  Any negative extension
already has twice the positive order, so must be the full group.
This proves that the table is exhaustive.

These are conjugacy classes, not merely abstract groups.  For the
cyclic positive actions, mark orbits of equal kind move into each
other in the quotient sphere while the poles are retained.  For
the Klein four cases the quotient is a mirror lune: its marked
corner and the numbers of points on its two mirror edges and in
its interior determine the decoration up to homeomorphism.  The
three distributions just listed exhaust its possibilities.  The
other actions are rigid as described above.  A homeomorphism of
a decorated quotient lifts to a conjugacy of its sphere action.
Conversely the permutation action on the five punctures together
with the orientation sign is invariant under mapping-class
conjugacy.  The twenty models give distinct conjugacy classes of
subgroups of $S_5\times C_2$, as checked exactly in
\texttt{check\_census.py}; hence none of the twenty classes can
merge in the mapping class group.  The induced finite mapping-class
action is faithful, since a positive element in its kernel fixes
all five marks, and a negative element cannot be isotopic to the
identity.  This finishes the classification.
\end{proof}

\statement{E.37}{Proved}{The nine vertex families contain 542 labeled objects}
The vertex structures displayed in cote 148, scan 31, have the
following cardinalities and are each one transitive $S_5$-set:
\[
\begin{array}{c|r|r}
 \text{type}&\text{stabilizer order}&\text{cardinality}\\\hline
 I_{10}&10&12\\
 I_6&6&20\\
 I_4&4&30\\
 II_2&2&60\\
 II'_2&2&60\\
 II_1&1&120\\
 III_2&2&60\\
 III_1&1&120\\
 III'_2&2&60
\end{array}
\]
Their disjoint union has $542$ objects in nine $S_5$-orbits.
\begin{proof}
The first three counts can be read without a geometric drawing:
there are $(5-1)!/2=12$ unoriented pentagons; there are
$\binom53\,2=20$ choices of a triangle with its orientation
torsor identified with the remaining pair; and there are
$5\cdot(4-1)!=30$ distinguished-point oriented squares.
A pentagon with one distinguished vertex has $12\cdot5=60$
structures.  Equivalently choose a three-set $J$, a point $P\in J$,
and a bijection from the complementary two-set to $J\setminus P$:
$10\cdot3\cdot2=60$.  These are the parameters of $II_2,II'_2$
and $III_2$ in the diagram, with their respective geometric
realizations.  The fully pinned objects $II_1,III_1$ have trivial
stabilizer and therefore $5!=120$ labelings.  For $III'_2$, choose
$P$ in five ways and an oriented square with the extra codiagonal
choice on the remaining four points in twelve ways, giving $60$;
its stabilizer is the indicated double transposition.  Equivalently
all nine sets are the coset sets $S_5/H$ for the listed stabilizers.
Every set is transitive by relabeling, so orbit--stabilizer gives
the whole table.  Summing the three levels gives
$62+240+240=542$.  \texttt{check\_census.py} constructs and
checks these finite $S_5$-sets exactly.
\end{proof}

\statement{E.38}{Refuted}{The fourteen arrow families form fourteen $S_5$-orbits}
For the arrow families and local degrees displayed in cote 148,
scans 31 and 33, the sizes are
\[
 60,120;\quad60,120,60,60,120;\quad
 120,120,120,120,360,120,60.
\]
They sum to $1620$.  They cannot form fourteen $S_5$-orbits.
The size-$360$ family from $III_1$ to $II_1$ has three free
$S_5$-orbits, and the indicated remaining families give thirteen
orbits.  The resulting count is sixteen, with sizes
$5\cdot60+11\cdot120$.
\begin{proof}
The two $f$-families contribute $60+120=180$ arrows.  The five
middle families contribute $420$.  The seven lower families,
including the family with three arrows out of each of the
$120$ fully pinned $III_1$-objects, contribute $1020$.
This gives $180+420+1020=1620$.

An orbit of a group of order $120$ cannot have $360$ elements.
More precisely, stabilizers of arrows are contained in the
stabilizers of their sources.  Since $III_1$ is a free $S_5$-set,
every arrow with this source has trivial stabilizer.  Thus its
$360$ arrows split into exactly three orbits of size $120$.
For the remaining families, the indicated auxiliary datum is
transitive under the stabilizer of the source: it is a vertex of
a triangle, one of two equivalent choices, or a uniquely determined
choice, as appropriate.  Their arrow sets are correspondingly
$S_5/C_2$ or $S_5$, with the sizes listed above.  Therefore the
fourteen \emph{families} give $13+3=16$ orbits.

The computation file records each source, target, cardinality,
source degree and target degree, and constructs the corresponding
coset-orbit model.  This arithmetic does not establish relations
among the geometric arrows.  Direct inspection of
scan 33 confirms the total $1620$ and the $3\times120$ entry,
although overwritten intermediate products are inconsistent and hard
to read.  Crucially, the word \emph{orbites} following $14$ is
struck out in the original.  Thus the summary in R9, which presents
fourteen orbits as a clean assertion, is too strong.  Fourteen is
the correct number of displayed arrow families, and sixteen is the
correct number of $S_5$-orbits for that diagram.  We do not attribute
the uncancelled fourteen-orbit assertion to the manuscript.
\end{proof}

\statement{E.39}{Refuted}{The undecorated square pyramid has no real structure}
If the phrase on cote 148, scan 31, excluding $I_4$ from the real
structures is interpreted as a statement about the undecorated
five-point complex configuration, it is false.  For
$K=\{\infty,1,i,-1,-i\}$, complex conjugation preserves $K$,
and its full extended automorphism group is $D_4$ of order eight.
\begin{proof}
Conjugation fixes $\infty,1,-1$, exchanges $i,-i$, and normalizes
the positive generator $z\mapsto iz$ by inversion.  The group
it generates with that rotation is $D_4$.  Since E.30 gives
positive order four, there are no further negative elements.
Thus $G=\widetilde G=C_4$ is not correct for the undecorated
configuration when $\widetilde G$ means all holomorphic and
antiholomorphic automorphisms.  An additional oriented decoration
can exclude the reflection, in which case the statement concerns
a different object.  The scan-31 diagram uses oriented square
data, whereas scan 36 explicitly includes the square pyramid
among real configurations.  We retain this distinction rather
than identify two different categories.
\end{proof}


\statement{E.40}{Proved}{The stable-tree counts on scan 36}
For five labeled marks, the stable genus-zero dual trees with zero,
one, and two edges have respectively $1,10,15$ combinatorial types.
The corresponding decorated counts written at the bottom of cote 148,
scan 36, are $12$ in the pentagonal top type, $60$ and $120$ for the
two indicated one-edge decorations, and $60$ for the indicated
two-edge decoration.  The additional pyramid and bipyramid counts
are $30$ and $20$.
\begin{proof}
A one-edge stable tree splits the labels as $2+3$, giving
$\binom52=10$.  A two-edge tree has a single central mark and two
unordered pairs at its ends: $5\cdot3=15$.  There are no further
stable tree shapes for five marks.  The stabilizers of a fixed
one-edge and two-edge labeled partition have orders $12$ and $8$.
The decorated objects of E.37 have respective stabilizers $2,1$
above the one-edge type and $2$ above the indicated two-edge type.
Thus their fibers contain $12/2=6$, $12/1=12$, and $8/2=4$
objects; multiplication gives $60,120,60$.  The other three counts
are the first three entries of E.37.  This verifies the displayed
numerical statements without identifying the distinct decorated
families having the same cardinality.
\end{proof}
